\chapter{Progress}\label{progress}

% Detailed statement of progress to date, including:
% - a summary of key research activities and outcomes
% - a list of any publications generated from your research (if applicable)
% - a summary of any major changes in research direction
% - a summary of thesis progress, chapters and their preliminary titles
% Anticipated progress in next 12 months including:
% - a brief timetable of planned research activities
% - a statement of the resources required to achieve these outcomes

The ODSA pipeline feeds on the REMA surface elevation model~\cite{REMA} ice thickness and the local ice-flow direction. The flow direction comes from the REMA surface gradient with a baseline of $10\times$ the ice thickness following the smoothing function recommendation from McCormack et al.~\cite{McCormack_2019}. ODSA also uses the $450~\mathrm{m}$ MEaSUREs v2 product~\cite{Rignot_2017}. At each point it gives the speed and the angle between the MEaSUREs velocity direction and the REMA flow direction. Allowing us to obtain the velocity uncertainties per window. A $50~\mathrm{km}$ sliding window at 50\% overlap is the unit in the spectral analysis of the bed RES data, (Not every window is 50 km. A segment shorter than 50 km is analysed as a single window the length of the segment).

The detrended bed profile for each Hann tapered window is analysed with a Lomb--Scargle periodogram (over 500 geomspace frequencies), which accommodates uneven along-track sample spacing~\cite{VanderPlas_2018}. 


\begin{figure}
    \includegraphics[scale=0.37]{figures/catalogue_examples.png}
    \caption{} 
    \label{fig:catalogue_examples}
\end{figure}